\documentclass[../../main.tex]{subfiles}
\begin{document}
\begin{oframed}
  \subfile{problem}
\end{oframed}

{\bf [解答]}
題意から
\begin{align}
  0 < a < b < 2a \label{eq:1}
\end{align}
の元で考える．
長方形$R$の4頂点を$W(0,0)$，$X(0,a)$，$Y(b,a)$，$Z(b,0)$とする（$XY=b$，$XW=a$）．
円$A$は性質(P)から2辺$XY,XW$に接するとしてよく，半径を$x$とすると中心は$(x,a-x)$である．
円$A$が長方形$R$内に納まるので
\begin{align}
  0<x<\dfrac{a}{2} \label{eq:2}
\end{align}
が必要である．

\begin{figure}[htb]
  \centering
  \begin{tikzpicture}[scale=1.0, >=stealth, every node/.style={font=\small}]

    % --- 1. パラメータ設定 (a = 4.6, b = 7.0, x = 1.6) ---
    \def\a{4.6}
    \def\b{7.0}
    \def\x{1.6}

    % 各円の半径計算（解説の公式値）
    % r1 = (-\sqrt{x} + \sqrt{a})^2 \approx 0.774
    \pgfmathsetmacro{\rOne}{( -sqrt(\x) + sqrt(\a) )^2}
    % r2 = a + b - \sqrt{2ab} - x \approx 1.975
    \pgfmathsetmacro{\rTwo}{\a + \b - sqrt(2*\a*\b) - \x}
    % r3 = (-\sqrt{x} + \sqrt{b})^2 \approx 1.907
    \pgfmathsetmacro{\rThree}{( -sqrt(\x) + sqrt(\b) )^2}
    % r4 = (3 - 2\sqrt{2})x \approx 0.275
    \pgfmathsetmacro{\rFour}{(3.0 - 2.0*sqrt(2.0))*\x}

    % --- 2. 各頂点・中心の座標定義 ---
    \coordinate (W) at (0, 0);
    \coordinate (Z) at (\b, 0);
    \coordinate (Y) at (\b, \a);
    \coordinate (X) at (0, \a);

    \coordinate (CenterA)  at (\x, {\a - \x});
    \coordinate (CenterC1) at (\rOne, \rOne);
    \coordinate (CenterC2) at ({\b - \rTwo}, \rTwo);
    \coordinate (CenterC3) at ({\b - \rThree}, {\a - \rThree});
    \coordinate (CenterC4) at (\rFour, {\a - \rFour});

    % --- 3. 補助線（軸投影・中心間） ---
    \draw[densely dotted, thick, gray!70!black] (CenterA) -- (\x, 0) node[below=1pt, font=\footnotesize, black] {$x$};
    \draw[densely dotted, thick, gray!70!black] (CenterA) -- (0, {\a - \x}) node[left=2pt, font=\footnotesize, black] {$a-x$};

    \draw[densely dotted, thick, gray!70!black] (CenterA) -- (CenterC1);
    \draw[densely dotted, thick, gray!70!black] (CenterA) -- (CenterC2);
    \draw[densely dotted, thick, gray!70!black] (CenterA) -- (CenterC3);
    \draw[densely dotted, thick, gray!70!black] (CenterA) -- (CenterC4);

    % --- 4. 長方形 WXYZ の描画 ---
    \draw[line width=1.2pt, gray!80!black] (W) -- (Z) -- (Y) -- (X) -- cycle;

    \node[below left=2pt] at (W) {$W(0,0)$};
    \node[below right=2pt] at (Z) {$Z(b,0)$};
    \node[above right=2pt] at (Y) {$Y(b,a)$};
    \node[above left=2pt] at (X) {$X(0,a)$};

    % --- 5. 円 A の描画 ---
    \draw[line width=1.2pt, blue!85!black] (CenterA) circle (\x);
    \node[circle, fill=blue!85!black, inner sep=1.3pt] at (CenterA) {};
    \node[blue!85!black, font=\footnotesize] at ($(CenterA) + (0.55, 0.48)$) {$A(x,\, a-x)$};

    % --- 6. 外接する 4 円 C_1, C_2, C_3, C_4 の描画と中心座標 ---
    % [C_1]: 中心 (r_1, r_1)
    \draw[thick, red!85!black] (CenterC1) circle (\rOne);
    \node[circle, fill=red!85!black, inner sep=1.2pt] at (CenterC1) {};
    \node[right, red!85!black, font=\footnotesize] (lblC1) at (1.85, 0.55) {$C_1(r_1,\, r_1)$};
    \draw[->, thin, red!85!black] (lblC1.west) to[out=180, in=-15] ($(CenterC1) + (0.08, -0.02)$);

    % [C_2]: 中心 (b - r_2, r_2)
    \draw[thick, red!85!black] (CenterC2) circle (\rTwo);
    \node[circle, fill=red!85!black, inner sep=1.2pt] at (CenterC2) {};
    \node[right, red!85!black, font=\footnotesize] (lblC2) at ({\b + 0.45}, 1.6) {$C_2(b-r_2,\, r_2)$};
    \draw[->, thin, red!85!black] (lblC2.west) to[out=180, in=10] ($(CenterC2) + (0.08, 0.01)$);

    % [C_3]: 中心 (b - r_3, a - r_3)
    \draw[thick, red!85!black] (CenterC3) circle (\rThree);
    \node[circle, fill=red!85!black, inner sep=1.2pt] at (CenterC3) {};
    \node[above, red!85!black, font=\footnotesize] (lblC3) at (4.9, {\a + 0.45}) {$C_3(b-r_3,\, a-r_3)$};
    \draw[->, thin, red!85!black] (lblC3.south) to[out=-90, in=80] ($(CenterC3) + (0.02, 0.08)$);

    % [C_4]: 中心 (r_4, a - r_4)
    \draw[thick, red!85!black] (CenterC4) circle (\rFour);
    \node[circle, fill=red!85!black, inner sep=1.1pt] at (CenterC4) {};
    \node[above right, red!85!black, font=\footnotesize] (lblC4) at (0.75, {\a + 0.45}) {$C_4(r_4,\, a-r_4)$};
    \draw[->, thin, red!85!black] (lblC4.south west) to[out=200, in=60] ($(CenterC4) + (0.05, 0.05)$);

    % --- 7. 寸法線（長方形の幅 b, 高さ a） ---
    \draw[<->, gray!70!black] (0, -0.60) -- (\b, -0.60) node[midway, below=1pt, black] {$b$};
    \draw[<->, gray!70!black] (-1.05, 0) -- (-1.05, \a) node[midway, left=1pt, black] {$a$};

  \end{tikzpicture}
  \caption{長方形 $WXYZ$ と円 $A$，および外接円 $C_k$ の中心座標}
\end{figure}

\paragraph{(1)}

円$A$に外接する性質(P)の円は，対称性から次の4つで尽くされる：
\begin{itemize}
  \item 頂点$W$で2辺$XW,WZ$に接する円$C_1$（半径$r_1$，中心$(r_1,r_1)$）
  \item 頂点$Z$で2辺$WZ,ZY$に接する円$C_2$（半径$r_2$，中心$(b-r_2,r_2)$）
  \item 頂点$Y$で2辺$ZY,YX$に接する円$C_3$（半径$r_3$，中心$(b-r_3,a-r_3)$）
  \item 頂点$X$で円$A$と同じ2辺$XY,XW$に接し，$A$より角$X$寄りにある円$C_4$（半径$r_4$，中心$(r_4,a-r_4)$）
\end{itemize}

このうち$C_4$は自明に存在するから，残りの$C_k\, (k=1,2,3)$の存在条件について調べれば良い．
これらの円は長方形の内部に納まるから，\cref{eq:2}と同様の条件
\begin{align}
  0 < r_k < \dfrac{a}{2} \label{eq:3}
\end{align}
を満たす．

さらにそれぞれの円が$A$と外接するから，円の中心間の距離がちょうど円の半径の和に等しいから
\begin{align}
  (x-r_1)^2+(a-x-r_1)^2=(x+r_1)^2 \label{eq:4}\\
  (b-x-r_2)^2+(a-r_2-x)^2=(x+r_2)^2 \label{eq:5}\\
  (b-x-r_3)^2+(x-r_3)^2=(x+r_3)^2 \label{eq:6}
\end{align}
である．

\begin{figure}[htb]
  \centering
  \begin{tikzpicture}[scale=1.1, >=stealth, every node/.style={font=\small}]

    % --- 1. パラメータ設定 (a = 4.6, b = 7.0, x = 1.6, r2 = 1.975) ---
    \def\a{4.6}
    \def\b{7.0}
    \def\x{1.6}
    \pgfmathsetmacro{\rTwo}{\a + \b - sqrt(2.0*\a*\b) - \x}

    % 頂点座標
    \coordinate (W) at (0, 0);
    \coordinate (Z) at (\b, 0);
    \coordinate (Y) at (\b, \a);
    \coordinate (X) at (0, \a);

    % 円の中心座標
    \coordinate (CenterA) at (\x, {\a - \x});
    \coordinate (CenterC2) at ({\b - \rTwo}, \rTwo);

    % --- 2. 座標軸と長方形 WXYZ ---
    \draw[->, thick, gray!80!black] (-0.7, 0) -- (\b+0.8, 0) node[right, black] {$x$};
    \draw[->, thick, gray!80!black] (0, -0.7) -- (0, \a+0.8) node[above, black] {$y$};

    \node[below left=2pt, gray!80!black] at (W) {$W(0,0)$};
    \node[below right=2pt, gray!80!black] at (Z) {$Z(b,0)$};
    \node[above right=2pt, gray!80!black] at (Y) {$Y(b,a)$};
    \node[above left=2pt, gray!80!black] at (X) {$X(0,a)$};

    % 長さ b, a の補助寸法線
    \draw[<->, gray!70!black] (0, -0.82) -- (\b, -0.82) node[midway, below=1pt, black] {$b$};
    \draw[<->, gray!70!black] (-1.05, 0) -- (-1.05, \a) node[midway, left=1pt, black] {$a$};

    % 長方形の外枠
    \draw[thick, gray!80!black] (W) -- (Z) -- (Y) -- (X) -- cycle;

    % --- 3. 座標軸への投影（点線と目盛り） ---
    \draw[densely dotted, thick, blue!70!black] (CenterA) -- (\x, 0) node[below=1pt, blue!85!black, font=\footnotesize] {$x$};
    \draw[densely dotted, thick, blue!70!black] (CenterA) -- (0, {\a - \x}) node[left=2pt, blue!85!black, font=\footnotesize] {$a-x$};

    \draw[densely dotted, thick, red!70!black] (CenterC2) -- ({\b - \rTwo}, 0) node[below=1pt, red!85!black, font=\footnotesize] {$b-r_2$};
    \draw[densely dotted, thick, red!70!black] (CenterC2) -- (0, \rTwo) node[left=2pt, red!85!black, font=\footnotesize] {$r_2$};

    % --- 4. 円 A と C_2，および中心間距離 ---
    \draw[line width=1.2pt, blue!85!black] (CenterA) circle (\x);
    \draw[line width=1.2pt, red!85!black] (CenterC2) circle (\rTwo);

    \draw[very thick, <->, magenta!80!black] (CenterA) -- (CenterC2);
    \node[circle, fill=blue!85!black, inner sep=1.5pt] at (CenterA) {};
    \node[circle, fill=red!85!black, inner sep=1.5pt] at (CenterC2) {};

    \node[blue!85!black, font=\Large] at ({\x + 0.5}, {\a - \x - 0.5}) {$A$};
    \node[red!85!black, font=\Large] at ({\b - \rTwo - 0.6}, {\rTwo + 0.6}) {$C_2$};
    \node[magenta!80!black, font=\large, fill=white, inner sep=2pt]
    at ($(CenterA)!0.5!(CenterC2)$) {$x + r_2$};

  \end{tikzpicture}
  \caption{長方形 $WXYZ$ 内の円 $A$ と $C_2$ の中心座標（軸投影）と中心間距離 $x+r_2$（例）}
\end{figure}

まず\cref{eq:4}で左辺第一項と右辺を展開して
\begin{align}
  (a-x-r_1)^2 = 4xr_1 \\
  a-x-r_1 = 2\sqrt{xr_1} \\
  \sqrt{r_1} = -\sqrt{x} \pm \sqrt{a}
\end{align}
であるが，$\sqrt{r_1}>0$より複号正が正しく，
\begin{align}
  \sqrt{r_1} = -\sqrt{x} + \sqrt{a} \label{eq:7}
\end{align}
である．全く同様に\cref{eq:6}も解けて
\begin{align}
  \sqrt{r_3} = -\sqrt{x} + \sqrt{b} \label{eq:8}
\end{align}
である．

最後に\cref{eq:5}は$x+r_2$についてまとめると
\begin{align}
  & (x+r_2)^2 -2(a+b)(x+r_2)+\left(a^2+b^2\right) = 0 \\
  \therefore
  x+r_2 & = (a+b) \pm \sqrt{(a+b)^2-(a^2+b^2)} \\
  & = (a+b) \pm \sqrt{2ab}
\end{align}
である．\cref{eq:2,eq:3}より$0<x+r_2<a$だから複号負が正しく，
\begin{align}
  x+r_2
  = a+b-\sqrt{2ab} \label{eq:9}
\end{align}
となる．以上\cref{eq:7,eq:8,eq:9}をまとめて
\begin{align}
  \begin{dcases}
    \sqrt{r_1} = -\sqrt{x} + \sqrt{a} \\
    r_2 = a+b-\sqrt{2ab} - x \\
    \sqrt{r_3} = -\sqrt{x} + \sqrt{b}
  \end{dcases} \label{eq:10}
\end{align}
である．これが\cref{eq:3}を満たすような$x$の範囲を求めれば良いから\cref{eq:3}に代入して
\begin{align}
  &
  \begin{dcases}
    0 < -\sqrt{x} + \sqrt{a} < \sqrt{\dfrac{a}{2}} \\
    0 < a+b-\sqrt{2ab} - x < \dfrac{a}{2} \\
    0 < -\sqrt{x} + \sqrt{b} < \sqrt{\dfrac{a}{2}}
  \end{dcases} \\
  \therefore
  &
  \begin{dcases}
    \sqrt{a} - \sqrt{\dfrac{a}{2}} < \sqrt{x} <  \sqrt{a}  \\
    \left(\sqrt{b}-\sqrt{\dfrac{a}{2}}\right)^2 <  x < a+b-\sqrt{2ab} \\
    \sqrt{b} - \sqrt{\dfrac{a}{2}} < \sqrt{x} < \sqrt{b}
  \end{dcases}
\end{align}
だが，\cref{eq:1,eq:2}より厳しい方の条件を採用して
\begin{align}
  \left(\sqrt{b} - \sqrt{\dfrac{a}{2}}\right)^2 < x < \dfrac{a}{2} \label{eq:11}
\end{align}
である．

\paragraph{(2)}

まずは$C_k$の内でどれが$B$に該当するかを見るため，半径$r_k$の比較を行う．
(1)で計算しなかった$r_4$について$r_1,r_2,r_3$と同様に考えると
\begin{align}
  &(x+r_4)^2 = (x-r_4)^2 + (x-r_4)^2 = 2(x-r_4)^2\\
  & r_4^2-6xr_4 + x^2 = 0 \\
  \therefore
  & r_4 = \left(3\pm 2\sqrt{2}\right) x
\end{align}
を得る．ここで\cref{eq:1,eq:11}より複号正の方は
\begin{align}
  \left(3+ 2\sqrt{2}\right) x
  & > \left(3+ 2\sqrt{2}\right)\left(\sqrt{b} - \sqrt{\dfrac{a}{2}}\right)^2 \quad (\because\text{\cref{eq:11}}) \\
  & > \left(3+ 2\sqrt{2}\right)\left(1 - \dfrac{\sqrt{2}}{2}\right)^2 a \quad (\because\text{\cref{eq:1}}) \\
  & = \left(3+ 2\sqrt{2}\right)\dfrac{3-2\sqrt{2}}{2} a  \\
  & = \dfrac{a}{2}
\end{align}
となって\cref{eq:3}を満たさず不適だから複号負の方が正しく，
\begin{align}
  r_4 = \left(3 - 2\sqrt{2}\right) x \label{eq:12}
\end{align}
である．以上\cref{eq:7,eq:8,eq:9,eq:12}をまとめると
\begin{align}
  \begin{dcases}
    r_1 = \left(-\sqrt{x} + \sqrt{a}\right)^2 \\
    r_2 = a+b-\sqrt{2ab} - x \\
    r_3 = \left(-\sqrt{x} + \sqrt{b}\right)^2 \\
    r_4 = \left(3 - 2\sqrt{2}\right) x
  \end{dcases} \label{eq:13}
\end{align}
だから，以下これらの大小を比較しよう．

まず，$r_1,r_4$について
\begin{align}
  r_1-r_4
  &= \left(-\sqrt{x} + \sqrt{a}\right)^2 - \left(3 - 2\sqrt{2}\right) x \\
  &= 2\left(-1+\sqrt{2}\right) x - 2\sqrt{ax} + a \\
  &= 2\left(-1+\sqrt{2}\right)\left(\sqrt{x} - \dfrac{\sqrt{a}}{2\left(-1+\sqrt{2}\right)}\right)^2 + a - \dfrac{a}{2\left(-1+\sqrt{2}\right)} \\
  &= 2\left(-1+\sqrt{2}\right)\left(\sqrt{x} - \dfrac{(1+\sqrt{2})\sqrt{a}}{2}\right)^2 + a - \dfrac{(1+\sqrt{2})a}{2} \\
  &= 2\left(-1+\sqrt{2}\right)\left(\sqrt{x} - \dfrac{(1+\sqrt{2})\sqrt{a}}{2}\right)^2 + \dfrac{(1-\sqrt{2})a}{2}
\end{align}
よりこれを$x$の関数とみなすと\cref{eq:2}の範囲では単調減少で，端点$x=\dfrac{a}{2}$で値が$0$だから区間内でこの値も正である．
従って
\begin{align}
  r_1 > r_4 \label{eq:14}
\end{align}
である．

次に$r_2$と$r_3$は
\begin{align}
  r_3 - r_2
  &= \left(-\sqrt{x} + \sqrt{b}\right)^2 - \left(a+b-\sqrt{2ab} - x\right) \\
  &= 2x - 2\sqrt{bx} +\sqrt{2ab} - a \\
  &= 2\left(\sqrt{x}-\dfrac{\sqrt{2a}}{2}\right)\left(\sqrt{x}-\dfrac{2\sqrt{b}-\sqrt{2a}}{2}\right) \\
  &< 0
\end{align}
より
\begin{align}
  r_3 < r_2 \label{eq:15}
\end{align}
である．

ついで$r_3$と$r_1$は
\begin{align}
  r_3-r_1
  &=\left(-\sqrt{x} + \sqrt{b}\right)^2 -  \left(-\sqrt{x} + \sqrt{a}\right)^2 \\
  &= -2\sqrt{bx} + 2\sqrt{ax} + b - a \\
  &= -2\sqrt{x}(\sqrt{b}-\sqrt{a}) + (\sqrt{b}-\sqrt{a})(\sqrt{b}+\sqrt{a}) \\
  &= \left[\sqrt{b}+\sqrt{a} - 2\sqrt{x}\right](\sqrt{b}-\sqrt{a})
\end{align}
となるが，これは$x$について単調減少で\cref{eq:11}より$x\to\dfrac{a}{2}$で最小となり，
\begin{align}
  r_3-r_1
  & > \left[\sqrt{b}+\sqrt{a} - 2\sqrt{\dfrac{a}{2}}\right](\sqrt{b}-\sqrt{a}) \\
  & = \left[\sqrt{b}+(1-\sqrt{2})\sqrt{a}\right](\sqrt{b}-\sqrt{a}) \\
  & > \left[\sqrt{a}+(1-\sqrt{2})\sqrt{a}\right](\sqrt{b}-\sqrt{a}) \quad (\because\text{\cref{eq:1}})\\
  & > 0
\end{align}
より
\begin{align}
  r_1 < r_3 \label{eq:16}
\end{align}
である．

以上\cref{eq:14,eq:15,eq:16}より
\begin{align}
  r_4 < r_1 < r_3 < r_2 \label{eq:17}
\end{align}
となり，二番目に半径が大きいのは$C_3$だから$B=C_3$である．

よって円$A$と円$B$の面積の和$S(x)$は\cref{eq:13}より
\begin{align}
  S(x)
  &= \pi x^2 + \pi r_3^2 \\
  &= \pi \left[x^2 + \left(-\sqrt{x} + \sqrt{b}\right)^4 \right] \label{eq:18}
\end{align}
である．$X=\sqrt{x}$とおくと
\begin{align}
  S(x)
  &= \pi\left[X^4 + \left(X - \sqrt{b}\right)^4 \right]
\end{align}
だから微分して
\begin{align}
  \dfrac{1}{\pi}\dfrac{dS}{dX}
  &= 4X^3 + 4\left(X-\sqrt{b}\right)^3
\end{align}
ここで$s^3+t^3 = (s+t)(s^2-ts+t^2)$の因数分解を用いて
\begin{align}
  \dfrac{1}{\pi}\dfrac{dS}{dX}
  &= 4\left(2X-\sqrt{b}\right)\left(X^2-X(X-\sqrt{b})+(X-\sqrt{b})^2\right) \\
  &= 4\left(2X-\sqrt{b}\right)\left(X^2-\sqrt{b}X+b\right)
\end{align}
である．第二項は正だから$S(x)$の増減表は以下のようになる．

\begin{table}[htb]
  \centering
  \caption{$S(x)$の増減表}
  \begin{tabular}{c|ccccc}
    $x$ & $\left(\sqrt{b}-\sqrt{a/2}\right)^2$ & & $b/4$ & & $a/2$ \\
    \hline
    $X$ & $\sqrt{b}-\sqrt{a/2}$ & & $\sqrt{b}/2$ & & $\sqrt{a/2}$ \\
    \hline
    $S'$ & & $-$ & $0$ & $+$ & \\
    \hline
    $S$ & & $\searrow$ & & $\nearrow$ &
  \end{tabular}
\end{table}

よって$X=\dfrac{\sqrt{b}}{2}$の時，すなわち$x=\dfrac{b}{4}$で$S$は最小となり，
この時\cref{eq:13}より
\begin{align}
  r_3
  &=\left(\sqrt{b}-\sqrt{\dfrac{b}{4}}\right)^2 \\
  &=\left(\dfrac{\sqrt{b}}{2}\right)^2=\dfrac{b}{4}
\end{align}
だから\cref{eq:18}に代入して求める面積の最小値は
\begin{align}
  \min S=\pi\left\{\left(\dfrac{b}{4}\right)^2+\left(\dfrac{b}{4}\right)^2\right\}=\dfrac{b^2}{8}\pi
\end{align}

\newpage

\end{document}
